\documentclass[11pt]{article}

\usepackage[margin=1in]{geometry}
\usepackage{amsmath,amssymb}
\usepackage{array}
\usepackage{enumitem}
\usepackage{hyperref}
\hypersetup{hidelinks}
\emergencystretch=2em

\title{Scaling Book Exercises: Section 3}
\author{Lucas Yan}
\date{}

\newcommand{\AG}{\operatorname{AllGather}}
\newcommand{\AR}{\operatorname{AllReduce}}
\newcommand{\RS}{\operatorname{ReduceScatter}}
\newcommand{\ATA}{\operatorname{AllToAll}}

\begin{document}
\maketitle

\noindent
\textbf{Chapter:} Sharded Matrices and How to Multiply Them\\
\textbf{Source scan:} \texttt{../handwritten/section\_03\_handwritten.pdf}

\paragraph{Notation.}
$W$ denotes bidirectional ICI bandwidth unless stated otherwise, $C$ denotes
per-accelerator FLOPs/s, and $X,Y,Z$ denote mesh-axis sizes as well as their
axis names when the meaning is clear. Arrays are in bfloat16 unless specified,
so each element occupies two bytes.

\section*{Exercise 1}

The mesh is $\{X:4,Y:8,Z:2\}$ and $A[I_X,J,K,\ldots]$ is sharded only over
$X$. Each device stores one quarter of $A$, while the data is replicated over
$Y$ and $Z$. Across all 64 devices,
\[
64\frac{|A|}{4}=16|A|.
\]
Thus the ratio of total bytes to one logical copy is
\[
\boxed{16}.
\]

\section*{Exercise 2}

For $B=1024$, $D=4096$, mesh $\{X:4,Y:4,Z:4\}$, and TPU v4p bandwidth
$W=9.0\times10^{10}$ bytes/s:

\begin{enumerate}[label=(\alph*)]
\item Gathering $[B_X,D_Y]$ over $X$ moves the $Y$-sharded fraction of the
array:
\[
T_{\AG_X}=\frac{2BD}{YW}
=\frac{2(1024)(4096)}{4(9.0\times10^{10})}
\approx\boxed{23\ \mu\mathrm{s}}.
\]

\item Gathering over both $X$ and $Y$ moves the full array while using two
mesh axes:
\[
T_{\AG_{XY}}=\frac{2BD}{2W}
\approx\boxed{46\ \mu\mathrm{s}}.
\]

\item An AllReduce is a ReduceScatter followed by an AllGather. The local
output shard has $2BD/(XY)$ bytes, so
\[
T_{\AR_Z}=\frac{4BD}{XYW}
\approx\boxed{11.6\ \mu\mathrm{s}}.
\]
\end{enumerate}

\section*{Exercise 3}

$[B_X]$ contains only $2B=256$ bytes globally and 64 bytes per device. Its
bandwidth time is negligible, but a bidirectional gather around an axis of size
four needs two hops. At roughly $1\ \mu$s per hop,
\[
\boxed{T\approx2\ \mu\mathrm{s}}.
\]

\section*{Exercise 4}

For $X[B,D]\,Y[D_X,F]\to Z[B,F]$, compare two strategies.

\paragraph{Strategy 1: gather weights, then multiply.}
\[
T_1=\max\left(\frac{2BDF}{C},\frac{2DF}{W}\right).
\]
It performs $2BDF$ FLOPs per device and communicates $2DF$ bytes.

\paragraph{Strategy 2: local multiply, then AllReduce.}
\[
T_2=\max\left(\frac{2BDF}{XC},\frac{4BF}{W}\right).
\]
It performs $2BDF/X$ FLOPs per device and AllReduces a $2BF$-byte output.

Strategy 2 is compute-bound only if
\[
X<\frac{DW}{2C}.
\]
With TPU v5p values $C/W\approx2550$ and $D\approx8000$, this requires
$X<1.6$, so Strategy 2 is normally communication-bound. Strategy 1 becomes
compute-bound at $B>C/W\approx2550$. If $B<2550$, Strategy 2 wins when
$D>2B$; for large $B$, it wins when $D>2C/W\approx5100$. Hence the
local-matmul/AllReduce strategy is usually better for large hidden dimensions.

\section*{Exercise 5}

On a $4\times4\times4$ mesh, let $P=XYZ=64$. To minimize latency while
returning a replicated result, shard one non-contracting output dimension,
perform the local matmul, then AllGather the output. Two symmetric choices are
\[
A[I_{XYZ},J]B[J,K]\to C[I_{XYZ},K]\xrightarrow{\AG_{XYZ}}C[I,K]
\]
or
\[
A[I,J]B[J,K_{XYZ}]\to C[I,K_{XYZ}]\xrightarrow{\AG_{XYZ}}C[I,K].
\]
For either choice,
\[
T=\max\left(\frac{2IJK}{64C},\frac{2IK}{3W}\right).
\]
When compute-bound, the calculation is 64 times faster than full replication.
In the communication-bound case, it still beats full replication when
\[
J>\frac{C}{3W}.
\]
Using the scan's $C=2.75\times10^{14}$ and $W=9.0\times10^{10}$ gives a
crossover near $J\approx1019$.

\section*{Exercise 6}

Consider a TPU v5e $4\times4$ mesh.

\begin{enumerate}[label=(\alph*)]
\item For
\[
A[I_X,J_Y]B[J_Y,K]\to C[I_X,K],
\]
both contracting axes are sharded identically. Local matmuls produce partial
sums unreduced over $Y$, followed by $\AR_Y$:
\[
T_{\mathrm{math}}=\frac{2IJK}{XYC},
\qquad
T_{\mathrm{comm}}=\frac{4IK}{XW}.
\]

\item For
\[
A[I_X,J]B[J_X,K_Y]\to C[I_X,K_Y],
\]
the shared use of mesh axis $X$ conflicts. First $\AG_X$ the $J_X$ dimension
of $B$, then perform the local matmul:
\[
T_{\mathrm{math}}=\frac{2IJK}{XYC},
\qquad
T_{\mathrm{comm}}=\frac{2JK}{YW}.
\]

\item For
\[
A[I_X,J]B[J,K_Y]\to C[I_X,K_Y],
\]
neither contracting dimension is sharded and the two non-contracting
dimensions use different mesh axes. No communication is required:
\[
T_{\mathrm{math}}=\frac{2IJK}{XYC},
\qquad T_{\mathrm{comm}}=0.
\]
\end{enumerate}

\section*{Exercise 7}

$W_{\mathrm{in}}$ and $W_{\mathrm{out}}$ each occupy
\[
2DF=2(8192)(32768)\approx537\text{ MB},
\]
so they must be sharded across the four-device $2\times2$ slice. A tensor-
parallel layout is
\[
\mathrm{In}[B,D_{XY}]\,W_{\mathrm{in}}[D,F_{XY}]
\to H[B,F_{XY}],
\]
\[
H[B,F_{XY}]\,W_{\mathrm{out}}[F_{XY},D]
\to \mathrm{Out}[B,D]\{U_{XY}\}
\xrightarrow{\RS_{XY,D}}\mathrm{Out}[B,D_{XY}].
\]
This uses an initial $\AG_{XY}$ of the input and a final $\RS_{XY}$ of the
output. The two matmuls perform
\[
\frac{4BDF}{XY}
\]
FLOPs per device, while total communication time is
\[
T_{\mathrm{comm}}=\frac{2BD}{W}.
\]
With $B=128$, $D=8192$, $F=32768$, $XY=4$,
$C=1.97\times10^{14}$ FLOPs/s, and $W=9.0\times10^{10}$ bytes/s,
\[
T_{\mathrm{math}}\approx174\ \mu\mathrm{s},
\qquad
T_{\mathrm{comm}}\approx23.3\ \mu\mathrm{s}.
\]
With overlap, the estimate is therefore $\boxed{174\ \mu\mathrm{s}}$; without
overlap, approximately $197\ \mu\mathrm{s}$.

\section*{Exercise 8}

The scan records the JAX collective mappings:
\begin{align*}
\AG &: \texttt{jax.lax.all\_gather},\\
\AR &: \texttt{jax.lax.psum},\\
\RS &: \texttt{jax.lax.psum\_scatter},\\
\ATA &: \texttt{jax.lax.all\_to\_all}.
\end{align*}
Each benchmark should be mapped over a named axis. Synchronize every result
with
\[
\texttt{jax.block\_until\_ready(result)}
\]
so dispatch latency is not mistaken for collective runtime. The handwritten
work does not include measured timing results, so this exercise remains an
implementation benchmark rather than a numerically verified result.

\section*{Exercise 9}

Let device $r$ own contracting-index set $S_r$. It computes the partial result
\[
C_r[n,k]=\sum_{m\in S_r}A[n,m]B[m,k].
\]
An AllReduce forms the replicated product:
\[
C[n,k]=\sum_r C_r[n,k]
=\sum_{m=0}^{M-1}A[n,m]B[m,k].
\]
If a sharded output is acceptable, replace the AllReduce with a ReduceScatter
along either $N$ or $K$. Gathering $A[N,M]$ communicates an amount proportional
to $NM$; reduce-scattering $C[N,K]$ communicates an amount proportional to
$NK$. Therefore,
\[
\frac{\text{gather-first cost}}{\text{matmul-first cost}}=\boxed{\frac{M}{K}}.
\]
Gather first when $M<K$, reduce-scatter after the matmul when $M>K$, and either
when $M=K$.

\section*{Exercise 10}

Let $D$ devices form a ring and let the logical matrix contain $N^2$ scalars.

\begin{enumerate}[label=(\alph*)]
\item On a unidirectional ring, AllGather or ReduceScatter sends a strip of
size $N^2/D$ for $D-1$ steps. Per link,
\[
V_{\mathrm{AG,uni}}=N^2\frac{D-1}{D}\longrightarrow N^2.
\]

\item For AllToAll, destination blocks shrink along the route:
\[
V_{\mathrm{ATA,uni}}
=\frac{N^2}{D^2}\sum_{j=1}^{D-1}j
=\frac{N^2(D-1)}{2D}
\longrightarrow\frac{N^2}{2}.
\]

\item Thus unidirectional AllToAll costs half as much as AllGather.

\item Bidirectionality halves AllGather time because each strip travels in
both directions.

\item Bidirectionality reduces AllToAll time by a factor of four: blocks choose
the shorter direction and travel at most $\lfloor D/2\rfloor$ hops.

\item Combining the ratios,
\[
\frac{T_{\mathrm{AG,bi}}}{T_{\mathrm{ATA,bi}}}
=\frac{(1/2)T_{\mathrm{AG,uni}}}{(1/4)T_{\mathrm{ATA,uni}}}
=\boxed{4}.
\]
\end{enumerate}

\end{document}
